% Generated TikZ assets for main_es.tex. These are original vector schematics.

\newcommand{\NSMiniSignal}[3]{%
  \draw[#1, line width=#2, domain=0:3.7, samples=100]
    plot (\x,{#3 + 0.18*sin(260*\x) + 0.08*sin(610*\x) + 0.05*cos(90*\x)});
}

\newcommand{\NSCoverVisual}{%
\begin{tikzpicture}[x=1cm,y=1cm,>=Latex]
  \fill[NSGray] (-0.4,-0.35) rectangle (16.8,7.2);
  \draw[NSBlue!18, line width=0.4pt] (0,0) grid[xstep=0.5,ystep=0.5] (16.4,6.8);
  \begin{scope}[shift={(8.2,3.75)}, scale=1.18]
    \draw[NSBlue, line width=1.15pt]
      (-2.1,0.15) .. controls (-2.45,1.25) and (-1.45,2.0) .. (-0.35,1.72)
      .. controls (0.05,2.25) and (1.15,2.05) .. (1.45,1.35)
      .. controls (2.28,1.05) and (2.15,-0.35) .. (1.42,-0.65)
      .. controls (1.18,-1.45) and (0.2,-1.62) .. (-0.25,-1.15)
      .. controls (-1.1,-1.65) and (-2.25,-0.95) .. (-2.1,0.15);
    \foreach \p/\q in {
      -1.55/0.75,-0.75/1.25,0.1/1.05,0.92/0.72,1.35/-0.1,
      0.55/-0.72,-0.38/-0.54,-1.2/-0.25,-1.72/0.05,-0.25/0.35}
      \fill[NSTeal!85!black] (\p,\q) circle (2.1pt);
    \draw[NSTeal!70!black, line width=0.8pt] (-1.55,0.75)--(-0.75,1.25)--(0.1,1.05)--(0.92,0.72)--(1.35,-0.1);
    \draw[NSTeal!70!black, line width=0.8pt] (-1.2,-0.25)--(-0.38,-0.54)--(0.55,-0.72)--(1.35,-0.1);
    \draw[NSTeal!70!black, line width=0.8pt] (-1.72,0.05)--(-1.2,-0.25)--(-0.25,0.35)--(0.1,1.05);
    \draw[NSOrange!80!black, line width=1.1pt] (-0.75,1.25) to[bend left=18] (0.55,-0.72);
    \draw[NSPurple!80!black, line width=1.1pt] (-1.55,0.75) to[bend right=14] (0.92,0.72);
  \end{scope}
  \foreach \y/\c in {1.05/NSPetrol,0.72/NSBlue,0.39/NSOrange} {
    \draw[\c!85!black, line width=0.8pt, domain=0.7:15.8, samples=240]
      plot (\x,{\y + 0.09*sin(170*\x) + 0.035*sin(490*\x)});
  }
  \node[anchor=west, align=left, text=NSBlue, font=\bfseries\Large] at (0.55,6.42) {EEG de reposo \& conectividad funcional};
  \node[anchor=west, align=left, text=NSInk, font=\small] at (0.58,5.95) {Cuaderno técnico interno: BIDS, QC, CSD, wPLI, surrogates y FDR};
\end{tikzpicture}%
}

\newcommand{\NSStudyMap}{%
\begin{tikzpicture}[x=1cm,y=1cm,>=Latex]
  \node[nsnode, fill=NSBlue!6, minimum width=2.4cm] (ms1) at (0,2.2) {Pacientes EM\\T1};
  \node[nsnode, fill=NSTeal!7, minimum width=2.4cm] (ctrl) at (0,0.8) {Controles\\T1};
  \node[nsnode, fill=NSPurple!7, minimum width=2.4cm] (ms2) at (0,-0.6) {Pacientes EM\\T2};
  \node[nswarn, minimum width=2.4cm] (loss) at (0,-2.0) {Pérdidas de\\seguimiento};
  \node[nsnode, minimum width=1.25cm] (eo) at (3.0,1.25) {EO};
  \node[nsnode, minimum width=1.25cm] (ec) at (3.0,-0.15) {EC};
  \node[nsphase] (raw) at (5.15,0.55) {EEG raw\\BrainVision};
  \node[nsphase] (bids) at (7.0,0.55) {BIDS\\+ metadatos};
  \node[nsphase] (prep) at (8.95,0.55) {Filtrado\\ICA\\epochs};
  \node[nsphase] (conn) at (11.0,0.55) {CSD\\+ wPLI};
  \node[nsphase] (inf) at (13.05,0.55) {Surrogates\\+ FDR};
  \node[nsnode, fill=NSOrange!7, minimum width=2.1cm] (out) at (15.15,0.55) {Matrices\\tablas\\QC};
  \foreach \a in {ms1,ctrl,ms2} {
    \draw[nsarr] (\a.east) -- (eo.west);
    \draw[nsarr] (\a.east) -- (ec.west);
  }
  \draw[nssoftarr] (ms1.south) -- (ms2.north);
  \draw[nsarr] (eo.east) -- (raw.west);
  \draw[nsarr] (ec.east) -- (raw.west);
  \draw[nsarr] (raw) -- (bids);
  \draw[nsarr] (bids) -- (prep);
  \draw[nsarr] (prep) -- (conn);
  \draw[nsarr] (conn) -- (inf);
  \draw[nsarr] (inf) -- (out);
  \node[align=center, font=\scriptsize, text=NSInk] at (1.55,-2.0) {registradas\\como trazabilidad};
\end{tikzpicture}%
}

\newcommand{\NSExperimentalDesign}{%
\begin{tikzpicture}[x=1cm,y=1cm,>=Latex]
  \node[font=\bfseries\color{NSBlue}] at (2.7,3.0) {Comparación transversal};
  \node[font=\bfseries\color{NSBlue}] at (10.2,3.0) {Seguimiento longitudinal};
  \node[nsnode, fill=NSBlue!6] (em1) at (1.2,1.65) {EM T1};
  \node[nsnode, fill=NSTeal!7] (co) at (4.25,1.65) {Control T1};
  \node[nsnode, fill=NSBlue!6] (lt1) at (8.6,1.65) {EM T1};
  \node[nsnode, fill=NSPurple!7] (lt2) at (11.8,1.65) {EM T2};
  \draw[nsarr, <->] (em1) -- node[above,font=\scriptsize] {grupo} (co);
  \draw[nsarr, <->] (lt1) -- node[above,font=\scriptsize] {tiempo} (lt2);
  \foreach \x in {1.2,4.25,8.6,11.8} {
    \node[nsnode, minimum width=0.7cm] at (\x-0.45,0.35) {EO};
    \node[nsnode, minimum width=0.7cm] at (\x+0.45,0.35) {EC};
  }
  \draw[NSLine] (6.45,-0.15)--(6.45,3.35);
  \node[align=center,font=\scriptsize,text=NSInk] at (2.7,-0.52) {Contrasta patrones EM frente a controles\\sin asumir dirección causal};
  \node[align=center,font=\scriptsize,text=NSInk] at (10.2,-0.52) {Evalúa cambio intra-grupo\\cuando existe seguimiento T2};
\end{tikzpicture}%
}

\newcommand{\NSPipelineGlobal}{%
\begin{tikzpicture}[x=1cm,y=1cm,>=Latex]
  \foreach \i/\name/\y in {
    0/{IO\\BIDS}/0,
    1/{Filtrado}/0,
    2/{ICA}/0,
    3/{Segmentación}/0,
    4/{Baseline}/0,
    5/{Rechazo\\artefactos}/0,
    6/{CSD\\FFT}/0,
    7/{wPLI}/0,
    8/{Surrogates\\FDR}/0} {
    \node[nsphase, minimum width=1.55cm, minimum height=0.85cm] (p\i) at (1.6*\i,0) {\textbf{F\i}\\\name};
  }
  \foreach \i/\j in {0/1,1/2,2/3,3/4,4/5,5/6,6/7,7/8} \draw[nsarr] (p\i) -- (p\j);
  \draw[decorate, decoration={brace, amplitude=4pt}, NSBlue] (-0.75,-0.8)--(3.95,-0.8);
  \node[font=\scriptsize, text=NSBlue] at (1.6,-1.25) {armonización y limpieza};
  \draw[decorate, decoration={brace, amplitude=4pt}, NSPetrol] (4.05,-0.8)--(8.75,-0.8);
  \node[font=\scriptsize, text=NSPetrol] at (6.4,-1.25) {señal por segmentos};
  \draw[decorate, decoration={brace, amplitude=4pt}, NSOrange] (8.85,-0.8)--(13.55,-0.8);
  \node[font=\scriptsize, text=NSOrange!85!black] at (11.2,-1.25) {conectividad e inferencia};
\end{tikzpicture}%
}

\newcommand{\NSConnectivityPerspectives}{%
\begin{tikzpicture}[x=1cm,y=1cm,>=Latex]
  \node[font=\bfseries\Large,text=NSBlue] at (6.2,3.75) {Tres perspectivas sobre las mismas regiones};
  \node[font=\small,text=NSInk!70] at (6.2,3.35) {Estructural = soporte \quad Funcional = sincronía \quad Efectiva = dirección};
  \foreach \x/\title/\c/\caption in {
    1.85/Estructural/NSBlue/Vías físicas,
    6.2/Funcional/NSTeal/Señales correlacionadas,
    10.55/Efectiva/NSPurple/Influencia temporal} {
    \node[draw=\c!35, fill=\c!8, rounded corners=3pt, minimum width=3.4cm, minimum height=2.35cm] (panel\x) at (\x,1.55) {};
    \node[font=\bfseries,text=\c!80!black] at (\x,2.55) {\title};
    \coordinate (n1) at (\x-0.9,1.55);
    \coordinate (n2) at (\x,2.15);
    \coordinate (n3) at (\x+0.9,1.55);
    \coordinate (n4) at (\x,0.78);
    \foreach \p/\lab in {n1/1,n2/2,n3/3,n4/4} {
      \fill[\c!85!black] (\p) circle (4.2pt);
      \node[font=\tiny,text=NSInk] at ($(\p)+(0,-0.23)$) {\lab};
    }
    \node[font=\scriptsize,text=NSInk!80] at (\x,0.33) {\caption};
  }
  \draw[NSBlue!60,line width=1pt] (0.95,1.55)--(1.85,2.15)--(2.75,1.55)--(1.85,0.78)--cycle;
  \draw[NSBlue!60,line width=1pt] (1.85,2.15)--(1.85,0.78);
  \draw[NSTeal!65!black,line width=1pt] (5.3,1.55)--(6.2,2.15)--(7.1,1.55)--(6.2,0.78)--cycle;
  \draw[NSTeal!65!black,line width=1pt] (5.3,1.55)--(6.2,0.78);
  \draw[NSPurple!75!black,->,line width=1pt] (9.65,1.55)--(10.55,2.15);
  \draw[NSPurple!75!black,->,line width=1pt] (10.55,2.15)--(11.45,1.55);
  \draw[NSPurple!75!black,->,line width=1pt] (9.65,1.55)--(10.55,0.78);
  \draw[NSPurple!75!black,->,line width=1pt] (10.55,2.15)--(10.55,0.78);
  \draw[NSPurple!75!black,->,line width=1pt] (11.45,1.55)--(10.55,0.78);
\end{tikzpicture}%
}

\newcommand{\NSRawQCFlow}{%
\begin{tikzpicture}[x=1cm,y=1cm,>=Latex]
  \node[nsnode, fill=NSBlue!5, minimum width=1.8cm, minimum height=0.9cm] (raw) at (0,0) {Datos\\raw};
  \node[nsphase, minimum width=1.9cm] (metrics) at (2.5,0) {métricas\\RMS, rango};
  \node[nsphase, minimum width=1.9cm] (thr) at (5.1,0) {umbrales\\+ z-score};
  \node[nswarn, minimum width=2.3cm] (review) at (7.9,0) {inspección\\visual};
  \node[nsphase, minimum width=2.2cm] (decision) at (10.8,0) {decisión\\canal/época};
  \node[nsnode, fill=NSTeal!8, minimum width=1.7cm] (ready) at (13.3,0) {listo\\para filtrar};
  \foreach \a/\b in {raw/metrics,metrics/thr,thr/review,review/decision,decision/ready} \draw[nsarr] (\a)--(\b);
  \draw[NSLine, line width=0.8pt] (4.35,-1.15)--(6.0,-1.15);
  \foreach \x/\h/\c in {4.55/0.35/NSTeal,4.85/0.65/NSTeal,5.15/0.95/NSOrange,5.45/0.5/NSTeal,5.75/1.2/NSOrange} {
    \draw[\c!80!black,line width=3pt] (\x,-1.15)--(\x,{-1.15+\h});
  }
  \node[font=\scriptsize,text=NSInk,align=center] at (5.15,-1.65) {las métricas orientan,\\no sustituyen la revisión};
  \draw[NSOrange!75!black,line width=1.1pt] (8.85,-1.1) arc[start angle=200,end angle=-20,radius=0.62];
  \draw[NSOrange!75!black,-Latex,line width=1.1pt] (9.45,-1.1)--(9.78,-0.68);
  \node[font=\scriptsize,text=NSInk] at (8.85,-1.65) {umbral ajustable};
\end{tikzpicture}%
}

\newcommand{\NSFFTFlow}{%
\begin{tikzpicture}[x=1cm,y=1cm,>=Latex]
  \node[nsnode, fill=NSBlue!5, minimum width=1.9cm] (in) at (0,1.0) {$\mathbf{E}^{bl2}$\\segmento};
  \node[nsphase, minimum width=1.9cm] (dc) at (2.35,1.0) {DC\\$E-\mu$};
  \node[nsphase, minimum width=2.1cm] (win) at (4.75,1.0) {ventana\\$E\odot w$};
  \node[nsphase, minimum width=2.1cm] (pad) at (7.25,1.0) {zero-padding\\$N_{fft}$};
  \node[nsphase, minimum width=1.8cm] (fft) at (9.65,1.0) {rFFT\\$X(f)$};
  \node[nsnode, fill=NSOrange!7, minimum width=2.1cm] (pow) at (9.65,-0.35) {potencia\\$|X_\ell|^2$};
  \node[nsphase, minimum width=2.4cm] (corr) at (6.95,-0.35) {correcciones\\$\overline{w^2}$, plegado};
  \node[nsnode, fill=NSTeal!7, minimum width=2.4cm] (bands) at (4.0,-0.35) {promedio\\por banda};
  \foreach \a/\b in {in/dc,dc/win,win/pad,pad/fft,fft/pow,pow/corr,corr/bands} \draw[nsarr] (\a)--(\b);
  \node[font=\scriptsize,text=NSInk,align=center] at (1.7,-1.0) {$\delta,\theta,\alpha,\beta,\gamma$\\como resumen operativo};
\end{tikzpicture}%
}

\newcommand{\NSBIDSBrainVision}{%
\begin{tikzpicture}[x=1cm,y=1cm,>=Latex]
  \node[nsnode, fill=NSOrange!7] (vhdr) at (0,2.0) {\texttt{.vhdr}\\cabecera};
  \node[nsnode, fill=NSOrange!7] (vmrk) at (0,0.8) {\texttt{.vmrk}\\marcadores};
  \node[nsnode, fill=NSOrange!7] (eeg) at (0,-0.4) {\texttt{.eeg}\\señal binaria};
  \node[nsphase, minimum width=2.4cm] (convert) at (3.1,0.8) {renombrado\\+ enlaces internos};
  \node[nsnode, fill=NSBlue!5, minimum width=2.6cm] (bids) at (6.4,0.8) {\texttt{sub-*/ses-*/eeg/}\\BrainVision BIDS};
  \node[nsnode, fill=NSTeal!6] (part) at (9.4,2.0) {\texttt{participants.tsv}};
  \node[nsnode, fill=NSTeal!6] (elec) at (9.4,0.8) {\texttt{electrodes.tsv}};
  \node[nsnode, fill=NSTeal!6] (json) at (9.4,-0.4) {sidecars\\\texttt{.json}};
  \node[nsphase, fill=NSPurple!7, draw=NSPurple!70!black] (val) at (12.3,0.8) {BIDS\\Validator};
  \foreach \a in {vhdr,vmrk,eeg} \draw[nsarr] (\a.east) -- (convert.west);
  \draw[nsarr] (convert) -- (bids);
  \foreach \a in {part,elec,json} \draw[nsarr] (bids.east) -- (\a.west);
  \foreach \a in {part,elec,json} \draw[nsarr] (\a.east) -- (val.west);
\end{tikzpicture}%
}

\newcommand{\NSEEGMontage}{%
\begin{tikzpicture}[x=1cm,y=1cm,>=Latex]
  \draw[NSBlue, line width=1pt] (0,0) circle (2.65);
  \draw[NSBlue, line width=1pt] (-0.35,2.65)--(0,3.02)--(0.35,2.65);
  \draw[NSBlue, line width=0.8pt] (-2.65,0.55).. controls (-3.05,0.18) and (-3.05,-0.18)..(-2.65,-0.55);
  \draw[NSBlue, line width=0.8pt] (2.65,0.55).. controls (3.05,0.18) and (3.05,-0.18)..(2.65,-0.55);
  \foreach \name/\x/\y/\c in {
    Fpz/0/2.25/NSOrange,Fp2/0.75/2.1/NSOrange,
    F7/-2.0/1.35/NSOrange,F3/-0.95/1.35/NSOrange,Fz/0/1.45/NSOrange,F4/0.95/1.35/NSOrange,F8/2.0/1.35/NSOrange,
    FC5/-1.45/0.85/NSOrange,FC1/-0.55/0.78/NSOrange,FCz/0/0.9/NSPurple,FC2/0.55/0.78/NSOrange,FC6/1.45/0.85/NSOrange,
    FT9/-2.3/0.55/NSTeal,FT10/2.3/0.55/NSTeal,T7/-2.35/0/NSTeal,T8/2.35/0/NSTeal,TP9/-2.25/-0.75/NSTeal,TP10/2.25/-0.75/NSTeal,
    C3/-1.05/0.05/NSPetrol,Cz/0/0.05/NSPetrol,C4/1.05/0.05/NSPetrol,
    CP5/-1.45/-0.65/NSPetrol,CP1/-0.55/-0.65/NSPetrol,CP2/0.55/-0.65/NSPetrol,CP6/1.45/-0.65/NSPetrol,
    P7/-2.0/-1.35/NSPurple,P3/-0.95/-1.35/NSPurple,Pz/0/-1.45/NSPurple,P4/0.95/-1.35/NSPurple,P8/2.0/-1.35/NSPurple,
    O1/-0.85/-2.15/NSBlue,Oz/0/-2.25/NSBlue,O2/0.85/-2.15/NSBlue} {
    \fill[\c!80!black] (\x,\y) circle (2.2pt);
    \node[font=\tiny, text=NSInk] at (\x,\y+0.18) {\name};
  }
  \node[font=\scriptsize, align=left] at (5.1,1.6) {\textcolor{NSOrange}{\rule{8pt}{4pt}} frontal\\
    \textcolor{NSPetrol}{\rule{8pt}{4pt}} central\\
    \textcolor{NSTeal}{\rule{8pt}{4pt}} temporal\\
    \textcolor{NSPurple}{\rule{8pt}{4pt}} parietal\\
    \textcolor{NSBlue}{\rule{8pt}{4pt}} occipital};
  \node[font=\scriptsize, align=left, text=NSInk] at (5.1,-0.7) {31 EEG\\REF: FCz\\GND: Fpz};
\end{tikzpicture}%
}

\newcommand{\NSFilterCascade}{%
\begin{tikzpicture}[x=1cm,y=1cm,>=Latex]
  \node[nsnode, minimum width=2cm] (raw) at (0,1.3) {señal cruda};
  \node[nsphase] (n50) at (2.4,1.3) {notch\\50 Hz};
  \node[nsphase] (b100) at (4.8,1.3) {bandreject\\100 Hz};
  \node[nsphase] (hp) at (7.2,1.3) {high-pass\\0.5 Hz};
  \node[nsphase] (lp) at (9.6,1.3) {low-pass\\150 Hz};
  \node[nsnode, fill=NSTeal!6, minimum width=2cm] (clean) at (12.0,1.3) {señal\\filtrada};
  \foreach \a/\b in {raw/n50,n50/b100,b100/hp,hp/lp,lp/clean} \draw[nsarr] (\a)--(\b);
  \draw[NSBlue!85, line width=0.8pt, domain=-0.8:1.05, samples=90] plot ({\x+0.8},{-0.55+0.25*sin(520*\x)+0.12*sin(90*\x)});
  \draw[NSTeal!85!black, line width=0.8pt, domain=-0.8:1.05, samples=90] plot ({\x+11.2},{-0.55+0.18*sin(250*\x)});
  \draw[->, NSInk!70] (2.1,-1.0)--(10.0,-1.0) node[right,font=\tiny] {Hz};
  \draw[->, NSInk!70] (2.1,-1.0)--(2.1,0.05);
  \draw[NSPetrol, thick] (2.1,-0.15)--(2.35,-0.15)--(2.45,-0.85)--(2.55,-0.15)--(4.5,-0.15)--(4.6,-0.85)--(4.7,-0.15)--(8.9,-0.15)--(9.15,-0.95)--(10.0,-0.95);
  \node[font=\tiny] at (2.52,-1.22) {50};
  \node[font=\tiny] at (4.65,-1.22) {100};
  \node[font=\tiny] at (9.15,-1.22) {150};
\end{tikzpicture}%
}

\newcommand{\NSICAFlow}{%
\begin{tikzpicture}[x=1cm,y=1cm,>=Latex]
  \node[nsnode] (mix) at (0,0) {señales\\mezcladas};
  \node[nsphase] (ica) at (2.4,0) {FastICA};
  \node[nsnode, fill=NSBlue!5] (ic) at (4.9,0) {componentes\\independientes};
  \node[nswarn] (art) at (7.45,0.95) {ocular\\muscular\\línea\\saltos};
  \node[nsnode, fill=NSTeal!6] (keep) at (7.45,-0.75) {ICs\\retenidos};
  \node[nsphase] (rec) at (10.0,0) {reconstrucción};
  \node[nsnode, fill=NSTeal!7] (out) at (12.4,0) {EEG limpio};
  \foreach \a/\b in {mix/ica,ica/ic,ic/art,ic/keep,keep/rec,rec/out} \draw[nsarr] (\a)--(\b);
  \draw[nssoftarr] (art) -- node[right,font=\tiny] {excluir} (keep);
  \foreach \y/\c in {1.65/NSBlue,1.38/NSPetrol,1.11/NSOrange} \NSMiniSignal{\c}{0.55pt}{\y}
  \foreach \x in {4.2,4.65,5.1,5.55} \draw[NSPetrol, line width=0.5pt, domain=0:0.75, samples=40] plot ({\x+\x*0},{1.15+0.15*sin(500*\x)});
\end{tikzpicture}%
}

\newcommand{\NSEpochBaseline}{%
\begin{tikzpicture}[x=1cm,y=1cm,>=Latex]
  \node[nsnode] (cont) at (0,0.9) {señal\\continua};
  \node[nsphase] (ep) at (2.6,0.9) {epochs\\1 s};
  \node[nsphase] (bl) at (5.1,0.9) {baseline\\0--100 ms};
  \node[nsphase] (ar) at (7.7,0.9) {rechazo\\$\pm70\,\mu$V};
  \node[nsnode, fill=NSTeal!6] (qc) at (10.4,0.9) {tensor limpio\\+ QC};
  \foreach \a/\b in {cont/ep,ep/bl,bl/ar,ar/qc} \draw[nsarr] (\a)--(\b);
  \draw[NSBlue, line width=0.7pt, domain=-0.7:0.85, samples=120] plot ({\x+0.7},{-0.65+0.18*sin(420*\x)});
  \foreach \x in {2.05,2.45,2.85,3.25} \draw[NSPetrol, rounded corners=1pt] (\x,-0.95) rectangle ++(0.3,0.55);
  \fill[NSOrange!12] (4.35,-0.95) rectangle (4.8,-0.4);
  \draw[NSPetrol] (4.35,-0.68)--(5.85,-0.68);
  \draw[red!70!black,dashed] (7.0,-0.35)--(8.45,-0.35);
  \draw[red!70!black,dashed] (7.0,-1.0)--(8.45,-1.0);
  \draw[NSBlue, line width=0.65pt] (7.0,-0.7)..controls(7.35,-0.25)and(7.75,-1.22)..(8.45,-0.62);
\end{tikzpicture}%
}

\newcommand{\NSCSDWPLI}{%
\begin{tikzpicture}[x=1cm,y=1cm,>=Latex]
  \node[nsnode, fill=NSBlue!5] (sensor) at (0,1.0) {potenciales\\en sensor};
  \node[nsphase] (csd) at (2.9,1.0) {CSD\\Laplaciano};
  \node[nsnode, fill=NSTeal!6] (spec) at (5.8,1.0) {menor\\contaminación\\espacial};
  \node[nsphase] (wpli) at (8.9,1.0) {wPLI};
  \node[nsnode, fill=NSOrange!7] (mat) at (11.8,1.0) {matriz\\$C\times C$};
  \foreach \a/\b in {sensor/csd,csd/spec,spec/wpli,wpli/mat} \draw[nsarr] (\a)--(\b);
  \node[align=center, font=\small, text=NSBlue] at (5.9,-0.75)
    {$\displaystyle \mathrm{wPLI}_{ij}=\frac{\left|\mathbb{E}[\operatorname{Im}(z_i\overline{z_j})]\right|}{\mathbb{E}[|\operatorname{Im}(z_i\overline{z_j})|]+\epsilon}$};
  \draw[NSLine] (4.0,-0.2)--(7.8,-0.2);
  \node[font=\scriptsize, text=NSInk] at (5.9,-1.35) {conectividad funcional robusta frente a desfase cero};
\end{tikzpicture}%
}

\newcommand{\NSWPLIIntuition}{%
\begin{tikzpicture}[x=1cm,y=1cm,>=Latex]
  \node[font=\bfseries\color{NSBlue}] at (6.35,3.45) {Intuición operativa de wPLI};
  \node[nsnode, fill=NSBlue!4, minimum width=3.45cm, minimum height=2.25cm] (zero) at (1.95,1.85) {};
  \node[nsnode, fill=NSTeal!5, minimum width=3.45cm, minimum height=2.25cm] (lag) at (6.35,1.85) {};
  \node[nsnode, fill=NSOrange!6, minimum width=3.45cm, minimum height=2.25cm] (im) at (10.75,1.85) {};

  \node[font=\scriptsize\bfseries,text=NSBlue] at (1.95,2.82) {desfase cercano a $0^\circ$};
  \node[font=\scriptsize\bfseries,text=NSPetrol] at (6.35,2.82) {desfase no nulo estable};
  \node[font=\scriptsize\bfseries,text=NSOrange!85!black] at (10.75,2.82) {signo imaginario consistente};

  \begin{scope}[shift={(0.55,1.95)}]
    \draw[NSBlue, line width=0.85pt, domain=0:2.8, samples=100] plot (\x,{0.33*sin(240*\x)});
    \draw[NSTeal!80!black, line width=0.85pt, domain=0:2.8, samples=100] plot (\x,{-0.42+0.33*sin(240*\x+10)});
    \draw[NSInk!45,dashed] (1.4,-0.95)--(1.4,0.55);
  \end{scope}
  \node[font=\scriptsize,text=NSInk,align=center] at (1.95,0.86) {peso bajo\\si domina la mezcla común};

  \begin{scope}[shift={(4.95,1.95)}]
    \draw[NSBlue, line width=0.85pt, domain=0:2.8, samples=100] plot (\x,{0.33*sin(240*\x)});
    \draw[NSTeal!80!black, line width=0.85pt, domain=0:2.8, samples=100] plot (\x,{-0.42+0.33*sin(240*\x-65)});
    \draw[NSInk!45,dashed] (1.4,-0.95)--(1.4,0.55);
    \draw[<->,NSPetrol,thick] (1.08,0.43)--(1.4,0.43);
    \node[font=\tiny,text=NSPetrol] at (1.24,0.68) {$\Delta\phi$};
  \end{scope}
  \node[font=\scriptsize,text=NSInk,align=center] at (6.35,0.86) {peso alto\\si el desfase persiste};

  \begin{scope}[shift={(9.45,1.85)}]
    \draw[->,NSInk!75] (-0.15,0)--(2.55,0) node[right,font=\tiny] {Re};
    \draw[->,NSInk!75] (1.2,-0.95)--(1.2,0.95) node[above,font=\tiny] {Im};
    \foreach \x/\y in {0.45/0.52,0.82/0.35,1.05/0.62,1.35/0.43,1.65/0.66,1.95/0.38,2.15/0.55}
      \fill[NSBlue!65] (\x,\y) circle (1.8pt);
    \foreach \x/\y in {0.55/-0.45,0.9/-0.68,1.45/-0.42,1.75/-0.58,2.08/-0.35}
      \fill[NSOrange!75] (\x,\y) circle (1.8pt);
    \draw[NSPurple, line width=0.8pt, domain=0:2.45, samples=80]
      plot (\x,{-1.68+0.13*sin(360*\x)+0.04*sin(960*\x)});
  \end{scope}
  \node[font=\scriptsize,text=NSInk,align=center] at (10.75,0.82) {más muestras con el mismo signo\\aumentan el cociente};

  \node[nsphase, minimum width=3.7cm] (formula) at (3.95,-0.65)
    {$\displaystyle
    \mathrm{wPLI}_{ij}=
    \frac{\left|\mathbb{E}[\operatorname{Im}(z_i\overline{z_j})]\right|}
    {\mathbb{E}[|\operatorname{Im}(z_i\overline{z_j})|]+\epsilon}$};

  \begin{scope}[shift={(8.35,-1.45)}]
    \foreach \r in {0,...,3} {
      \foreach \c in {0,...,3} {
        \pgfmathsetmacro{\val}{18+14*\r+10*\c}
        \fill[NSPetrol!\val!NSOrange] (0.42*\c,0.42*\r) rectangle ++(0.39,0.39);
      }
    }
    \draw[NSInk!55] (0,0) rectangle (1.65,1.65);
    \node[font=\tiny,text=NSInk] at (0.82,-0.25) {matriz $C\times C$};
  \end{scope}

  \draw[nsarr] (zero.south) |- (formula.west);
  \draw[nsarr] (lag.south) -- (formula.north);
  \draw[nsarr] (im.south) |- (formula.east);
  \draw[nsarr] (formula.east) -- ++(1.8,0) |- (8.35,-0.62);
  \node[font=\scriptsize,text=NSInk,align=left] at (11.0,-0.65)
    {Salida por banda,\\condición y sujeto};
\end{tikzpicture}%
}

\newcommand{\NSSurrogatesFDR}{%
\begin{tikzpicture}[x=1cm,y=1cm,>=Latex]
  \node[nsnode] (orig) at (0,0.8) {señal\\original};
  \node[nsphase] (phase) at (2.4,0.8) {aleatorizar\\fase};
  \node[nsnode, fill=NSBlue!5] (null) at (4.9,0.8) {distribución\\nula};
  \node[nsphase] (pval) at (7.35,0.8) {$p$-valores\\empíricos};
  \node[nsphase] (fdr) at (9.8,0.8) {FDR\\BH};
  \node[nsnode, fill=NSTeal!7] (sig) at (12.25,0.8) {matriz\\significativa};
  \foreach \a/\b in {orig/phase,phase/null,null/pval,pval/fdr,fdr/sig} \draw[nsarr] (\a)--(\b);
  \draw[NSBlue, thick] (4.1,-0.9)--(6.0,-0.9);
  \foreach \x/\h in {4.2/0.25,4.45/0.6,4.7/0.95,4.95/1.15,5.2/0.85,5.45/0.45,5.7/0.18} \draw[NSPetrol, line width=2pt] (\x,-0.9)--(\x,{-0.9+\h});
  \draw[NSOrange, thick] (5.52,-0.9)--(5.52,0.45);
  \node[font=\tiny, text=NSOrange!85!black] at (5.72,0.55) {real};
  \node[align=center,font=\scriptsize,text=NSInk] at (6.35,-1.35) {$p_{ij}=(1+\#\{w^{(s)}_{ij}\ge w^{real}_{ij}\})/(1+N_s)$};
\end{tikzpicture}%
}

\newcommand{\NSOutputsDashboard}{%
\begin{tikzpicture}[x=1cm,y=1cm,>=Latex]
  \node[nsnode, fill=NSBlue!5, minimum width=2.2cm, minimum height=1.0cm] at (0,1.35) {logs\\por fase};
  \node[nsnode, fill=NSTeal!6, minimum width=2.2cm, minimum height=1.0cm] at (2.55,1.35) {figuras\\QC};
  \node[nsnode, fill=NSPetrol!6, minimum width=2.2cm, minimum height=1.0cm] at (5.1,1.35) {tablas\\resumen};
  \node[nsnode, fill=NSPurple!6, minimum width=2.2cm, minimum height=1.0cm] at (0,0) {snapshot\\config};
  \node[nsnode, fill=NSOrange!8, minimum width=2.2cm, minimum height=1.0cm] at (2.55,0) {hash\\commit};
  \node[nsnode, fill=NSGray, minimum width=2.2cm, minimum height=1.0cm] at (5.1,0) {carpeta\\por ejecución};
  \node[nsphase, minimum width=2.5cm, minimum height=2.35cm] (run) at (8.5,0.65) {\runid};
  \foreach \x/\y in {0/1.35,2.55/1.35,5.1/1.35,0/0,2.55/0,5.1/0} \draw[nsarr] (\x+1.1,\y) -- (run.west);
\end{tikzpicture}%
}
